\documentclass[10pt,letterpaper]{article}
\usepackage[letterpaper,margin=0.72in,headheight=20pt,footskip=28pt]{geometry}
\usepackage[T1]{fontenc}\usepackage{lmodern}\usepackage[table]{xcolor}\usepackage{array,longtable,booktabs,tabularx}\usepackage{enumitem}\usepackage[most]{tcolorbox}\usepackage{fancyhdr,lastpage}\usepackage{microtype}\usepackage{amsmath,amssymb}\usepackage[hidelinks,bookmarks=false]{hyperref}
\definecolor{navy}{HTML}{12304A}\definecolor{teal}{HTML}{137C8B}\definecolor{cuebg}{HTML}{EAF2F5}\definecolor{notebg}{HTML}{F8FAFB}\definecolor{line}{HTML}{B8C8D1}\definecolor{gold}{HTML}{C9922E}
\setlength{\parindent}{0pt}\setlength{\parskip}{4pt}\setlist[itemize]{leftmargin=1.25em,itemsep=1pt,topsep=2pt}\setlist[enumerate]{leftmargin=1.5em,itemsep=1pt,topsep=2pt}\renewcommand{\arraystretch}{1.18}
\pagestyle{fancy}\fancyhf{}\lhead{\small\color{navy}\textbf{CISSP CBK Cornell Notes}}\rhead{\small\color{teal}Domain 5: Identity and Access Management}\cfoot{\small\color{navy}Page \thepage\ of \pageref{LastPage}}\renewcommand{\headrulewidth}{0.4pt}\renewcommand{\headrule}{\hbox to\headwidth{\color{teal}\leaders\hrule height \headrulewidth\hfill}}
\newenvironment{cornell}{\rowcolors{2}{white}{notebg}\begin{longtable}{@{}>{\raggedright\arraybackslash\columncolor{cuebg}\color{navy}\bfseries}p{0.285\textwidth}>{\raggedright\arraybackslash}p{0.65\textwidth}@{}}\rowcolor{navy}\color{white}\textbf{Cue / Active Recall} & \color{white}\textbf{Detailed Notes}\\ \hline\endfirsthead\rowcolor{navy}\color{white}\textbf{Cue / Active Recall (cont.)} & \color{white}\textbf{Detailed Notes}\\ \hline\endhead}{\end{longtable}}
\newcommand{\noterow}[2]{#1 & #2\\\arrayrulecolor{line}\hline}\newcommand{\source}[1]{\par{\footnotesize\color{teal}\textit{Source: CBK, #1}}\par}
\newtcolorbox{summarybox}[1]{breakable,colback=cuebg,colframe=teal,boxrule=0.7pt,arc=2mm,title=\textbf{#1},fonttitle=\color{white},colbacktitle=teal}\newtcolorbox{exambox}[1]{breakable,colback=white,colframe=gold,boxrule=0.7pt,arc=2mm,title=\textbf{#1},fonttitle=\color{white},colbacktitle=gold}
\begin{document}\hypersetup{pageanchor=false}
\begin{titlepage}\pagecolor{navy}\color{white}\vspace*{1.0in}{\Large THE OFFICIAL (ISC)\textsuperscript{2} CISSP CBK REFERENCE\par}\vspace{0.25in}{\large Sixth Edition \textbullet\ 2021\par}\vfill{\Huge\bfseries Cornell Notes\par}\vspace{0.25in}{\LARGE Domain 5\par}\vspace{0.08in}{\LARGE\bfseries Identity and Access Management\par}\vfill{\large Arthur Deane and Aaron Kraus\par}\vspace{0.12in}{\normalsize Study and review companion\par}\vspace{0.35in}{\small Primary source: \textit{The Official (ISC)\textsuperscript{2} CISSP CBK Reference}, Sixth Edition, pp. 377--418.\par}\end{titlepage}
\hypersetup{pageanchor=true}\nopagecolor\color{black}

\section*{\color{navy}Domain Overview}
Identity and Access Management (IAM) establishes who or what is requesting access, proves the asserted identity, determines what that subject may do, and records activity so actions are attributable. Effective IAM controls both physical and logical resources and manages access from initial identity proofing through provisioning, use, review, role change, and deprovisioning.

\begin{summarybox}{Review Objectives}
Distinguish identification, authentication, authorization, and accountability; select authentication factors; explain federation and credential management; compare authorization models; manage the access lifecycle and privileged accounts; and explain OAuth/OIDC, SAML, Kerberos, RADIUS, and TACACS+.
\end{summarybox}

\section{\color{navy}Control Physical and Logical Access to Assets}
\source{pp. 378--386}
\begin{cornell}
\noterow{Four IAM questions}{\textbf{Identification}: who/what do you claim to be? \textbf{Authentication}: can you prove it? \textbf{Authorization}: are you allowed this action/resource now? \textbf{Accountability}: can actions be reliably attributed and reviewed?}
\noterow{Information access}{Apply classification, ownership, need-to-know, least privilege, and data-state protections. Access should be explicit, approved, reviewable, and revoked when the purpose ends; possession of a system account does not automatically justify every data permission.}
\noterow{System access}{Control interactive users, administrators, services, APIs, batch jobs, and machine identities. Use unique identities where accountability is required; protect shared/system accounts with compensating controls and managed credentials.}
\noterow{Device access}{Authenticate/enroll devices, establish ownership/trust/posture, protect local credentials and encryption keys, restrict removable/peripheral access, and revoke lost/noncompliant devices. Device trust should be evaluated separately from user identity.}
\noterow{Facility access}{Layer perimeter and internal zones using guards/reception, locks, badges, biometrics, mantraps where justified, visitor escort/logging, and monitoring. Physical authorization should follow role and be reviewed like logical access.}
\noterow{Application access}{Authenticate the subject, authorize each protected function/data request, maintain secure session state, and log sensitive actions. Do not rely on client-side UI hiding as authorization; enforce access on the trusted service side.}
\end{cornell}

\section{\color{navy}Identification and Authentication}
\source{pp. 387--401}
\begin{cornell}
\noterow{Identification vs. authentication}{Identification is an assertion such as username, device ID, or certificate subject. Authentication validates that assertion using an authenticator/factor. A username alone identifies; the password/token/biometric used afterward authenticates.}
\noterow{Identity management implementation}{Central IdM/IAM can support requests, approvals, provisioning, synchronization, password/credential services, federation/SSO, role/entitlement management, reporting, and deprovisioning. Centralization improves consistency but makes IAM infrastructure highly sensitive.}
\noterow{Type 1 factor}{\textbf{Something known}: password, PIN, passphrase, or knowledge answer. Protect against guessing, reuse, phishing, credential stuffing, insecure reset, and storage compromise; knowledge alone is not multiple factors.}
\noterow{Type 2 factor}{\textbf{Something possessed}: registered device, hardware token, smart card, certificate/private-key device, or badge. Protect issuance, binding, loss/replacement, cloning, and recovery processes.}
\noterow{Type 3 factor}{\textbf{Something the user is or does}: fingerprint, face/iris/hand geometry, voice, typing behavior, and other biometric characteristics. Biometrics are difficult to replace if compromised and constitute sensitive personal information.}
\noterow{Type 4 / contextual factor}{The edition also describes an emerging contextual category based on how/where a subject authenticates, such as device or location. It can drive adaptive requirements such as prompting MFA for a new device. Treat context as risk evidence, not a permanent trusted identity.}
\noterow{What makes MFA?}{Use factors from \textbf{different} categories. Password + PIN is still knowledge-only; password + authenticator device, or smart card + biometric, combines independent factor types and resists a single class of compromise.}
\noterow{Biometric FAR / FRR / CER}{\textbf{FAR}: false acceptance rate -- unauthorized subject accepted. \textbf{FRR}: false rejection rate -- authorized subject rejected. \textbf{CER}: crossover error rate where FAR and FRR are equal; lower CER generally indicates a more accurate biometric system. Tune sensitivity based on risk and usability.}
\noterow{Biometric design risks}{Address liveness/spoofing, privacy/PII storage, accessibility, cultural acceptance, sensor/environmental variation, bias, body changes, and fallback/recovery. A biometric template is not as easily rotated as a password.}
\noterow{Accountability}{Use unique identities, reliable authentication, authorization, logs/audit trails, timestamps, and protected records so actions can be linked to subjects/processes. Shared accounts weaken accountability unless tightly controlled and independently logged.}
\noterow{Session management}{After authentication, protect the session identifier/token: use strong unpredictable values, secure transport, appropriate cookie flags, limited lifetime/idle timeout, rotation after privilege change/login, logout/invalidation, and protection from fixation/hijacking/replay.}
\noterow{Registration and identity proofing}{Establish an identity before credential issuance. The assurance of evidence/verification should match risk. Validate authoritative information, prevent duplicate/fraudulent enrollment, securely deliver initial credentials, and keep recovery/re-proofing from becoming the weakest path.}
\noterow{Federated Identity Management (FIM)}{Establish trust so an identity authenticated in one domain can be accepted by another. Federation reduces separate credentials but creates dependency on identity-provider assurance, assertions/tokens, trust configuration, and timely revocation.}
\noterow{Credential management}{Bind authenticators to identities across issuance, activation, use, storage, renewal/rotation, recovery, suspension, revocation, and destruction. Protect secrets/keys and administrative workflows; credential reset is itself a high-risk authentication event.}
\noterow{Single Sign-On (SSO)}{Authenticate once and reuse trusted identity/session information across multiple services. SSO improves usability and central policy but increases the impact of IdP/session compromise, so strengthen the authentication, availability, monitoring, and recovery of the SSO service.}
\noterow{Just-in-Time (JIT)}{Create/provision access when a qualified user first needs a resource, based on role/policy rather than maintaining permanent accounts everywhere. JIT reduces standing administration and can shrink privilege exposure when paired with time/need constraints.}
\end{cornell}

\begin{exambox}{MFA Cue}
Two secrets are not two factors. MFA requires independent factor categories. Also separate \textbf{identity proofing} (establish the person/entity initially) from \textbf{authentication} (prove the established identity during access).
\end{exambox}

\section{\color{navy}Federated Identity with Third-Party Services}
\source{pp. 401--403}
\begin{cornell}
\noterow{On-premises federation}{An internal directory/IAM/SSO can extend assertions to external services through gateways/connectors. Protect the internal identity source, federation signing keys, trust metadata, connectors, and outbound claims.}
\noterow{Cloud IDaaS}{Identity as a Service can provide federation, MFA, lifecycle, and integration for many cloud/SaaS applications. Evaluate provider security, availability, breach impact, data/privacy location, administrative controls, logging, recovery, lockout/exit, and integration with authoritative HR/identity sources.}
\noterow{Hybrid IAM}{Combine internal identity systems with cloud federation/IDaaS. Define which source is authoritative, synchronize only needed attributes, prevent duplicate/stale identities, protect sync/federation channels, and ensure deprovisioning reaches both environments promptly.}
\end{cornell}

\section{\color{navy}Authorization Mechanisms}
\source{pp. 404--408}
\begin{cornell}
\noterow{Role-Based Access Control (RBAC)}{Permissions are assigned to roles representing job functions, and subjects inherit permissions from assigned roles. RBAC scales well and supports least privilege/separation of duties, but role design and role creep require governance.}
\noterow{Rule-Based Access Control (RuBAC)}{System applies predefined rules such as time, network, action, or system condition. Rules are centrally enforced and can be deterministic; do not confuse rule-based with role-based despite the similar acronym.}
\noterow{Mandatory Access Control (MAC)}{Central authority establishes security labels/classifications and clearances; users/owners cannot freely change access. Common in government/high-assurance settings where information flow follows hierarchical policy.}
\noterow{Discretionary Access Control (DAC)}{Object owner has discretion to grant/revoke access. Flexible and common in operating systems, but permissions can propagate or be granted inconsistently, increasing insider/misconfiguration risk.}
\noterow{Attribute-Based Access Control (ABAC)}{Policy evaluates attributes of subject, object/resource, action, and environment/context. More expressive and dynamic than simple roles, but correct attribute quality and policy complexity are key risks.}
\noterow{Risk-Based Access Control}{Dynamically changes requirements/decisions based on assessed risk such as device, location, behavior, transaction sensitivity, and anomalies. Can require step-up authentication or deny access when risk rises.}
\end{cornell}

\section{\color{navy}Identity and Access Provisioning Lifecycle}
\source{pp. 408--413}
\begin{cornell}
\noterow{Joiner-mover-leaver lifecycle}{Provision approved minimum access for new users; adjust access promptly when roles change; remove unnecessary/privileged access and credentials when employment/relationship ends. The goal is no stale entitlement surviving beyond business need.}
\noterow{Account access review}{Periodically recertify who has access and which permissions/accounts remain justified. Prioritize privileged, sensitive, dormant, orphaned, service, and high-risk entitlements; require accountable owners/managers to resolve exceptions.}
\noterow{Account usage review}{Review what user, service, and nonhuman accounts actually do, not only what they are entitled to do. Logs/analytics can reveal dormant accounts, unusual behavior, abuse, compromised credentials, or excessive privilege.}
\noterow{Provisioning / deprovisioning}{Provisioning creates/enables identity and access after approval; deprovisioning suspends/terminates/deletes credentials and sessions. Integrate with authoritative personnel/contract status and verify downstream systems actually removed access.}
\noterow{Role definition}{Create roles from business responsibilities and least privilege, map permissions to roles, assign accountable owners, test separation-of-duties conflicts, review memberships, and prevent uncontrolled role proliferation or permission accumulation.}
\noterow{Privilege escalation}{Elevation can be legitimate (approved administrative task) or malicious (exploit to obtain higher rights). Use privileged-access management, JIT/time limits, separate admin accounts, MFA, approval, session recording/logging, and rapid revocation.}
\noterow{Service/nonhuman accounts}{Service, application, API, automation, and managed accounts need owners, least privilege, managed secrets/keys, noninteractive restrictions when possible, rotation, monitoring, and lifecycle/decommission controls. Every nonhuman identity still requires accountability.}
\end{cornell}

\section{\color{navy}Authentication and Federation Systems}
\source{pp. 414--418}
\begin{cornell}
\noterow{OAuth}{Authorization framework: a resource owner grants a client delegated access to a protected resource through an authorization server, typically using tokens. OAuth is primarily about \textbf{authorization}, not proving user identity to the client.}
\noterow{OpenID Connect (OIDC)}{Authentication/identity layer built on OAuth 2.0. An OpenID Provider issues identity claims/tokens to a relying party, enabling web/mobile SSO-like authentication. Validate issuer, audience, signature, nonce/state, redirect URIs, and token lifetime according to implementation.}
\noterow{SAML}{XML-based federation framework for exchanging security assertions among a user agent, Identity Provider (IdP), and Service Provider (SP). The chapter identifies assertions, bindings, protocols, and profiles. Commonly used for enterprise browser SSO.}
\noterow{Kerberos realm}{Principals authenticate through a Key Distribution Center. The Authentication Server issues a Ticket Granting Ticket; the Ticket Granting Server issues service tickets that kerberized applications accept. Symmetric keys and time-sensitive tickets reduce repeated password exposure.}
\noterow{Kerberos risks}{Protect the KDC and privileged identities, synchronize time, secure endpoints/ticket caches, rotate service secrets, and monitor ticket abuse. A compromised central service or reusable ticket/key can have broad impact.}
\noterow{RADIUS}{Centralized Authentication, Authorization, and Accounting (AAA) protocol widely used for network/remote/802.1X access. It authenticates users and records session/accounting information; deploy with protected transport/shared-secret configuration appropriate to the environment.}
\noterow{TACACS+}{AAA protocol commonly used to administer network infrastructure. It centralizes device login/authorization and can control individual administrative commands, supporting separation and audit of network-management privileges.}
\end{cornell}

\section{\color{navy}Critical Comparisons}
\begin{center}\small
\begin{tabularx}{\textwidth}{>{\bfseries\color{navy}\raggedright\arraybackslash}p{0.22\textwidth}>{\raggedright\arraybackslash}X>{\raggedright\arraybackslash}X}
\toprule Concept & First idea & Distinguishing idea\\\midrule
Identify / authenticate & Assert an identity & Prove the assertion\\
Authenticate / authorize & Prove who/what & Decide permitted action/resource\\
RBAC / RuBAC & Permissions through job role & Predefined system rules/conditions\\
MAC / DAC & Central labels/clearances & Owner discretion over object permissions\\
OAuth / OIDC & Delegated authorization & Authentication/identity layer over OAuth 2.0\\
SAML / Kerberos & XML federation assertions & Ticket-based symmetric SSO in a realm\\
FAR / FRR & Unauthorized user accepted & Authorized user rejected\\
\bottomrule
\end{tabularx}
\end{center}

\section{\color{navy}Key Terms and Acronyms}
\begin{summarybox}{Rapid-Review Glossary}
\textbf{AAA} -- Authentication, Authorization, Accounting \quad
\textbf{ABAC} -- Attribute-Based Access Control \quad
\textbf{DAC} -- Discretionary Access Control \quad
\textbf{FAR/FRR/CER} -- biometric error metrics \quad
\textbf{FIM} -- Federated Identity Management \quad
\textbf{IAM} -- Identity and Access Management \quad
\textbf{IdP} -- Identity Provider \quad
\textbf{JIT} -- Just-in-Time \quad
\textbf{KDC} -- Key Distribution Center \quad
\textbf{MAC} -- Mandatory Access Control \quad
\textbf{OIDC} -- OpenID Connect \quad
\textbf{RBAC} -- Role-Based Access Control \quad
\textbf{SAML} -- Security Assertion Markup Language \quad
\textbf{SSO} -- Single Sign-On.
\end{summarybox}

\section{\color{navy}High-Yield Review}
\begin{itemize}
\item IAM sequence: establish identity, authenticate, authorize each access, preserve accountability, review, and deprovision.
\item Factors in this edition: Type 1 know, Type 2 have, Type 3 biometric is/does, plus emerging contextual Type 4.
\item MFA requires distinct factor categories; identity proofing/recovery must be as trustworthy as ordinary login.
\item RBAC, RuBAC, MAC, DAC, ABAC, and risk-based control answer different policy needs.
\item OAuth delegates authorization; OIDC adds authentication; SAML exchanges federation assertions; Kerberos uses tickets.
\item Lifecycle errors create orphaned/stale access; access and actual account usage both require review.
\end{itemize}

\section{\color{navy}Active-Recall Questions}
\begin{enumerate}
\item Distinguish identification, authentication, authorization, and accountability.
\item What are the three established authentication factor types in this edition, and what contextual Type 4 does it describe?
\item Why is password + PIN not MFA?
\item Define FAR, FRR, and CER.
\item How do owner-driven DAC and centrally labeled MAC differ?
\item Compare RBAC, rule-based access, ABAC, and risk-based access.
\item Why can SSO reduce password burden but increase systemic impact?
\item What is the difference between OAuth and OIDC?
\item Name the three primary SAML actors.
\item What are the AS, TGS, and service tickets in Kerberos?
\item Why must nonhuman accounts have owners and lifecycle controls?
\item What must happen when a user changes roles?
\end{enumerate}

\subsection*{\color{teal}Answer Key}
\begin{enumerate}
\item Identification asserts; authentication proves; authorization permits; accountability attributes actions.
\item Type 1 know, Type 2 have, Type 3 biometric is/does; Type 4 uses contextual evidence such as device/location to adapt authentication.
\item Both are knowledge factors; multiple credentials from one category are still single-factor.
\item FAR = unauthorized accepted; FRR = authorized rejected; CER = point at which the rates meet.
\item DAC lets an object owner grant access; MAC enforces centrally defined labels/clearances users cannot freely alter.
\item RBAC uses roles; RuBAC predefined rules; ABAC evaluates attributes; risk-based dynamically responds to assessed contextual risk.
\item It centralizes authentication, so compromise/outage of the SSO/IdP can affect many services.
\item OAuth delegates authorization; OIDC supplies authentication/identity claims on top of OAuth 2.0.
\item User agent, Identity Provider, Service Provider.
\item AS authenticates/helps issue initial ticket material; TGS issues service tickets after a TGT; service tickets are presented to target services.
\item Machine identities can hold powerful persistent permissions; ownership, secrets, review, and retirement preserve accountability.
\item Re-evaluate and promptly remove obsolete permissions while provisioning only the newly justified role access.
\end{enumerate}

\section{\color{navy}Cornell Summary}
\begin{summarybox}{Domain 5 Summary}
IAM is a lifecycle of trust decisions, not merely a login screen. Establish a valid identity, bind and protect authenticators, require stronger factors when risk warrants, authorize the least privilege for each request, federate identities through explicitly governed trust, record actions for accountability, review both entitlements and actual use, and remove access when roles or relationships end. Central services such as SSO/IdP/Kerberos simplify consistent control while simultaneously becoming high-value resilience and security dependencies.
\end{summarybox}

\section{\color{navy}Final Domain Checklist}
\begin{itemize}[label=\(\square\)]
\item I can control logical and physical access to information, systems, devices, facilities, and applications.
\item I can distinguish all identity/authentication concepts and factor/error metrics.
\item I can explain session, proofing, federation, credential, SSO, and JIT management.
\item I can compare RBAC, RuBAC, MAC, DAC, ABAC, and risk-based models.
\item I can manage joiner/mover/leaver, access/usage review, roles, and privileged/nonhuman accounts.
\item I can explain OAuth/OIDC, SAML, Kerberos, RADIUS, and TACACS+.
\end{itemize}
\vfill{\footnotesize\color{teal}Prepared as a study companion to Deane \& Kraus, \textit{The Official (ISC)\textsuperscript{2} CISSP CBK Reference}, 6th ed. Content is paraphrased and synthesized for study.}
\end{document}

